\chapter{Experiments and results}\label{ch:experiments}

This chapter presents the experiments as an argument rather than in
the order they ran: the core finding and its content--language
decomposition first (\S\ref{sec:exp_001}--\ref{sec:exp_001_debiased}),
then three attempts to kill it as a measurement artefact
(\S\ref{sec:exp_015}--\ref{sec:exp_016}), then its generalisation
across conflict families and the metrics that recover what cosine
averages away (\S\ref{sec:exp_011}--\ref{sec:exp_021}), and finally
the event-ladder corpus that varies time within a conflict
(\S\ref{sec:exp_007}).

\section{Cross-lingual ingroup pull on the Russo-Ukrainian
  case}\label{sec:exp_001}\label{sec:exp_002}\label{sec:exp_014}

Across the fourteen cosine-eligible providers spanning six
regulatory regimes, responses to the nine-question Russo-Ukrainian
prompt bank embed closer to the own-language Wikipedia anchor than
to the other-language anchors. Averaged over the matrix, the
row-centred diagonal pull is $+0.175$ on English, $+0.043$ on
Russian, and $+0.026$ on Ukrainian
(Table~\ref{tab:exp_014_headline};
Figure~\ref{fig:exp_014_forest}). The EN ${>}$ RU ${>}$ UK
ordering is unanimous at provider level: the English diagonal
exceeds the Russian, and the Russian the Ukrainian, in fourteen of
fourteen providers. Thirteen of the fourteen are positive on all
three diagonals; the exception is the locally hosted
TAIDE-Llama3-8B, whose Russian ($-0.004 \pm 0.012$) and Ukrainian
($-0.033 \pm 0.029$) diagonals sit at or below zero --- a first
sight of the small-open-weight scope condition developed in
\S\ref{sec:exp_016}. Under the two-level cluster bootstrap of
\S\ref{sec:stats} the matrix means carry 95\% intervals of
$[+0.153, +0.200]$ (EN), $[+0.016, +0.067]$ (RU), and
$[+0.003, +0.048]$ (UK); the EN advantage is decisively resolved
(EN$-$RU difference interval $[+0.10, +0.17]$), while the RU--UK
\emph{magnitude} difference is not ($[-0.02, +0.05]$). The
RU ${>}$ UK ordering is directionally consistent across all
fourteen providers but statistically unresolved at the matrix
level; because every provider answers the same nine questions,
provider-level unanimity adds little evidence beyond the
question-resampled interval, and the ordering is read as
suggestive rather than established. Because four of the nine
questions share the annexation anchor, the stricter resampling unit
is the anchor (six blocks) rather than the question; under an
anchor-block bootstrap the English interval is essentially
unchanged ($[+0.15, +0.19]$) while the Russian interval widens to
touch zero ($[-0.00, +0.07]$) and the Ukrainian stays marginally
positive ($[+0.00, +0.05]$) --- the contested-language raw pulls
are small enough that their resolution depends on what one treats
as the independent topical unit. The permutation test of \S\ref{sec:stats}
places the ingroup-pull association far outside its null: the
observed grand mean row-centred diagonal ($+0.081$) exceeds the
largest of $10{,}000$ label-shuffled replicates ($+0.019$;
$p < 10^{-4}$). On raw embeddings this null is partly trivial ---
shuffling response-language labels breaks the same-script pairing,
so shared script and lexicon alone push the observed diagonal above
it. The substantive version runs the identical permutation on the
language-axis-projected embeddings, where that trivial route is
removed: the debiased association remains far outside its null
(observed $+0.064$ against a null maximum of $+0.015$;
$p < 10^{-4}$).

\begin{table}[!htbp]
  \centering
  \resizebox{\textwidth}{!}{%
\begin{tabular}{lrrr}
    \hline
    Provider & EN pull & RU pull & UK pull \\
    \hline
    \texttt{c4ai-aya-expanse-32b}       & $+0.192 \pm 0.025$ & $+0.056 \pm 0.027$ & $+0.045 \pm 0.023$ \\
    \texttt{command-r7b-arabic-02-2025} & $+0.188 \pm 0.023$ & $+0.051 \pm 0.031$ & $+0.031 \pm 0.032$ \\
    \texttt{grok-3-mini}                & $+0.184 \pm 0.026$ & $+0.051 \pm 0.029$ & $+0.029 \pm 0.026$ \\
    \texttt{ernie-4.5-300b-a47b}        & $+0.182 \pm 0.025$ & $+0.050 \pm 0.034$ & $+0.034 \pm 0.019$ \\
    \texttt{claude-haiku-4-5}           & $+0.182 \pm 0.022$ & $+0.040 \pm 0.027$ & $+0.016 \pm 0.025$ \\
    \texttt{gpt-4o-mini}                & $+0.180 \pm 0.024$ & $+0.058 \pm 0.032$ & $+0.029 \pm 0.028$ \\
    \texttt{gemini-2.5-flash}           & $+0.180 \pm 0.029$ & $+0.048 \pm 0.032$ & $+0.040 \pm 0.017$ \\
    \texttt{glm-4.5}                    & $+0.175 \pm 0.030$ & $+0.048 \pm 0.032$ & $+0.039 \pm 0.020$ \\
    \texttt{ALLaM-7B} (local)           & $+0.173 \pm 0.023$ & $+0.031 \pm 0.015$ & $+0.002 \pm 0.031$ \\
    \texttt{qwen-plus}                  & $+0.172 \pm 0.024$ & $+0.046 \pm 0.032$ & $+0.036 \pm 0.022$ \\
    \texttt{jamba-mini-2-2026-01}       & $+0.171 \pm 0.021$ & $+0.040 \pm 0.025$ & $+0.033 \pm 0.029$ \\
    \texttt{deepseek-chat}              & $+0.168 \pm 0.023$ & $+0.045 \pm 0.028$ & $+0.031 \pm 0.015$ \\
    \texttt{TAIDE-Llama3-8B} (local)    & $+0.156 \pm 0.026$ & $-0.004 \pm 0.012$ & $-0.033 \pm 0.029$ \\
    \texttt{mercury-2}                  & $+0.149 \pm 0.042$ & $+0.037 \pm 0.031$ & $+0.033 \pm 0.018$ \\
    \hline
    Matrix mean                       & $+0.175$           & $+0.043$           & $+0.026$           \\
    \hline
  \end{tabular}}
  \caption[Per-provider ingroup pull on the Russo-Ukrainian
    case]{Row-centred diagonal ingroup pull per provider and
    language on the \texttt{ru\_uk\_core} sweep, ordered by EN pull.
    Intervals are 95\% normal-approximation CIs across the nine
    per-question row-centred diagonal means ($n{=}9$ topical
    clusters per cell), following the effective-sample-size policy
    of \S\ref{sec:stats}; the ten samples within a question are not
    treated as independent draws for topic-level inference. Source:
    \texttt{exp\_014\_qclustered\_cis.csv} (OpenAI panel; the
    four-embedder appendix table,
    Appendix~\ref{sec:appendix_exp_015_table}, reports within-cell
    dispersion on pooled cosines instead).}
  \label{tab:exp_014_headline}
\end{table}

\begin{figure}[!htbp]
  \centering
  \includegraphics[width=\linewidth]{exp_014_forest.png}
  \caption[Per-provider ingroup pull, forest view]{Forest view of
    Table~\ref{tab:exp_014_headline}: row-centred diagonal ingroup
    pull per provider and language, with question-clustered 95\%
    CIs ($n{=}9$ per cell), providers ordered
    by English pull. Dotted vertical lines mark the per-language
    matrix means. The EN ${>}$ RU ${>}$ UK ordering is unanimous
    across all fourteen providers; the only cells at or below zero
    are TAIDE's Russian and Ukrainian diagonals, with ALLaM's
    Ukrainian cell just above zero.}
  \label{fig:exp_014_forest}
\end{figure}

Provider identity moves the response geometry far less than
question or language do. Defining agreement as the mean raw cosine
between two providers' samples on the same (question, language)
cell, averaged over the 27 cells, cross-provider agreement among
the twelve frontier-scale providers spans $0.65$--$0.85$ (mean
$0.77$), approaching those same providers' within-provider
self-agreement of $0.77$--$0.92$ (computed identically between a
provider's own ten samples, identity pairs excluded). The two
locally hosted 7B models fall outside both ranges (self-agreement
$0.55$ and $0.56$; every pair involving one of them sits at
$0.44$--$0.57$) --- the stylistic-fingerprint finding of
\S\ref{sec:exp_016}, visible already here. Self-agreement is a
temperature-dependent quantity --- all sweeps sample at
temperature $1.0$, which compresses it and thereby shrinks the gap
to cross-provider agreement --- so the comparison is tilted
\emph{toward} provider interchangeability rather than neutral, and
is read with that in mind. At frontier scale the
data support a one-phenomenon reading rather than twelve
provider-specific phenomena. Re-running
the five-provider pilot at twice the original token cap and
re-embedding gives a mean absolute change of under $0.003$ in the
row-centred cells across the $5\times 3\times 3$ cube --- the
headline is invariant to sample-length saturation (per-cell deltas
in Appendix~\ref{sec:appendix_token_budget_delta}).

A geometric check closes off one further alternative reading.
Because the diagonal is row-centred, its attainable maximum is set
by the anchors alone: a response identical to its own-language
anchor would score one minus that anchor row's mean, so mutually
similar anchors mechanically compress the attainable pull. The
anchor--anchor cosines are asymmetric across languages (EN--RU
$0.68$, EN--UK $0.60$, RU--UK $0.72$ under the primary embedder),
giving mechanical ceilings of $0.24$ (EN), $0.20$ (RU), and $0.23$
(UK). The geometry therefore contributes to the English advantage,
but it cannot produce the observed pattern: the ceiling spread
($0.01$--$0.04$) is a fraction of the observed spread
($0.13$--$0.15$); the ceilings place Ukrainian \emph{above}
Russian while the observed diagonals run the other way; and
ceiling-normalised diagonals preserve EN ${\gg}$ RU ${>}$ UK
($73\%$, $21\%$, and $11\%$ of ceiling attained). Anchor
geometry is a minor scaling factor on the headline, not its
source.

\section{Content versus language: the debiasing
  decomposition}\label{sec:exp_001_debiased}

Projecting onto the orthogonal complement of the language subspace
(\S\ref{sec:debiasing}) separates content-driven from
language-driven pull. Across the fourteen-provider matrix the
English diagonal collapses by $12\%$ ($+0.175 \to +0.153$), the
Russian diagonal by $28\%$ ($+0.043 \to +0.031$), and the Ukrainian
diagonal by $66\%$ ($+0.026 \to +0.009$), with the two locally
hosted 7B providers' debiased Ukrainian pull at or below zero.
Rerunning the two-level cluster bootstrap of \S\ref{sec:stats} on
the projected embeddings puts the debiased English diagonal at
$[+0.13, +0.18]$ --- decisively positive after the control --- the
debiased Russian residual at $[+0.005, +0.06]$ --- marginally
positive --- and the debiased Ukrainian residual at
$[-0.02, +0.03]$, which includes zero: after removing the language
axis, a genuine Ukrainian content pull can no longer be
distinguished from none at all, and the Russian residual is
resolved only barely (the response-fitted bracket below puts it at
$-0.001$). Refitting the language axis from
the response set's own per-language centroids --- the bracket
motivated in \S\ref{sec:debiasing}, which removes response-register
language variance the Wikipedia-fitted axis may miss --- pushes the
projection further in the same direction: EN $+0.101$, RU
$-0.001$, UK $-0.020$. Under either fit the qualitative reading is
the same and the second fit makes it sharper. The
EN ingroup pull is largely content-driven; most of the UK pull was
a lexical artefact of the shared script --- and the residual third
that survives the control-fitted projection is itself within noise,
a reading the
cross-embedder panel of \S\ref{sec:exp_015} further narrows.

A second view from the off-topic same-language baseline of
\S\ref{sec:topic_vs_lang} converges on the same reading. Topic lift
is strictly positive in every (event, language) cell and the
asymmetry across the three languages tracks the lexical-vs-content
cut.

\section{Methodology robustness}\label{sec:exp_015}\label{sec:exp_015_debiased}\label{sec:exp_017}

Re-embedding the full response and anchor corpus under three further
embedders --- Alibaba \verb|text-embedding-v3|, Google
\verb|gemini-embedding-001|, and Yandex \verb|text-search-doc| ---
gives a four-way panel. The \emph{sign} of cross-lingual ingroup
pull is method-robust: every per-language mean diagonal stays
positive under all four embedders (at per-provider granularity the
one exception is the same TAIDE cells already flagged in
Table~\ref{tab:exp_014_headline}). The \emph{magnitude ordering} EN ${\gg}$ RU ${>}$ UK is
not: Alibaba and Yandex flatten it to EN ${\approx}$ UK ${>}$ RU,
and Gemini compresses every diagonal toward zero with English only
marginally ahead. The anchor geometry of \S\ref{sec:exp_001}
accounts for much of this magnitude shrinkage --- embedders whose
anchors are mutually closer have lower mechanical ceilings
(Gemini's anchor--anchor cosines run $0.86$--$0.92$, capping its
diagonals near $0.07$--$0.09$) --- a further reason to compare
orderings rather than magnitudes across embedders.
The Russian-trained Yandex embedder gives Ukrainian responses
\emph{more} pull than Russian responses ($+0.090$ vs $+0.042$),
falsifying the naive ``Russian-corpus embedder would inflate RU''
rebuttal.

Cross-embedder debiasing complicates the
\S\ref{sec:exp_001_debiased} reading. Under
Alibaba and Yandex the Ukrainian pull \emph{survives} the
projection, so the lexical-vs-content split is itself a property of
the OpenAI embedding space rather than of the LLM responses.

A separate challenge comes from \citet{qi2023crosslingual}, who
report Russian--Ukrainian BLOOM-tokenizer overlap of $0.76$ ---
predicting the cluster fusion as a vocabulary artefact. Measuring
pairwise tokenizer overlap on the universal-control corpus for the
three inspectable embedder tokenizers gives scores well below the
BLOOM reference (\verb|cl100k| $0.53$, Qwen-proxy $0.42$, Cohere
$0.31$). Re-running the unsupervised analysis under the
lowest-overlap embedder redistributes which questions fuse but
preserves total fusion magnitude. Subword-vocabulary overlap is
ruled out as the dominant cause.

\section{Unsupervised structure recovers the experimental
  design}\label{sec:exp_016}\label{sec:exp_016_outliers}\label{sec:exp_016_anchor_projection}

Running EVoC \citep{mcinnes2024evoc} on the full response set
\emph{without} provider, question, or language labels produces
a resolution layer of twenty-nine clusters against the twenty-seven
(question $\times$ language) cells of the design. The reported
layer is simply the one nearest the design count --- a
presentational choice; the claim rests on the purity and
association values, not the cluster count. Per-cluster purity is
high on question ($\geq 0.91$) and language ($\geq 0.74$) and low
on provider ($0.13$--$0.19$); because raw purity is not comparable
across axes with different numbers of classes, the comparison is
scored by adjusted mutual information (\S\ref{sec:stats}), which
gives question $0.66$, language $0.36$, and provider $0.08$ ---
the provider axis carries several times less cluster information
than either design axis. The result holds across all four
embedders, whose chosen layers land near the design count
(19--35 clusters) with the same purity profile
(Appendix~\ref{sec:appendix_evoc_purity}).

Three provider outliers complicate the reading and each becomes its
own finding. YandexGPT collapses into a tight four-cluster pocket
rather than spreading across the (qid $\times$ language) cells.
TAIDE-Llama3-8B and ALLaM-7B sit at $59\%$ and $37\%$ on a
top-three-cluster concentration ranking against $12\%$--$24\%$ for
the other twelve providers. The scope condition this evidences is
operational rather than a parameter-count law: the two models that
carry stylistic fingerprints are the two \emph{locally hosted}
open-weight 7B models, while \verb|command-r7b| --- also a 7B
open-weight model, but served through its provider's API --- sits
inside the frontier band on every geometry measure here (and
resurfaces separately as the largest stance-swing outlier in
\S\ref{sec:exp_004}). ``Frontier-scale'' in this thesis therefore
names the twelve providers that behave homogeneously on the
embedding geometry, and provider-agnostic claims are scoped to that
empirically defined set rather than to a parameter threshold.

Projecting the Wikipedia anchors as points inside the EVoC fit
closes the loop with cosine: anchors fall into the cluster of their
matched-language responses, which is the structural form of
positive ingroup pull. Two disclosure notes bound this validation:
it is available for the four of nine questions whose anchors landed
in a real cluster (the q01 and q05 anchors fall into EVoC's noise
partition, and the three stance-slot questions share the annexation
anchor and split across stance-bound clusters), and the two metrics
share the embedding space even though they share no anchor
arithmetic.

The projection also exposes a structural fusion that the diagonal
cosine conceals. On Crimea, Maidan, and Bandera the \emph{Russian
and Ukrainian Wikipedia anchors land in a single joint cluster},
and both Russian and Ukrainian responses occupy it at
$82$--$89\%$ --- symmetrically, since it is one cluster --- while
every English anchor sits in a separate English-response cluster.
(On Crimea the Ukrainian anchor falls into the noise partition, so
the Ukrainian responses' $87\%$ co-clustering there is with the
Russian anchor by default.) The fusion is a property of how the
embedding space treats the Russian- and Ukrainian-language
treatments of these events jointly, not a directional claim that
Ukrainian responses prefer the Russian treatment over their own ---
the clustering cannot distinguish those readings when the two
anchors are inseparable. It is this \emph{fusion}, directional or
not, that the tokenizer-overlap control of \S\ref{sec:exp_017} was
built to challenge --- \citet{qi2023crosslingual} predict exactly
this pattern from subword-vocabulary overlap alone. That the
pattern survives at matched cluster resolution under the
lowest-overlap embedder is what licenses reading it as
content-level fusion rather than vocabulary leakage; the
directional Russian-to-Ukrainian reading rests on the debiasing
decomposition (\S\ref{sec:exp_001_debiased}), not on this
projection. Figure~\ref{fig:exp_016_anchor_map} shows
the projection.

\begin{figure}[!htbp]
  \centering
  \includegraphics[width=\linewidth]{exp_016_anchor_map.png}
  \caption[Anchor projection in the unsupervised embedding
    map]{UMAP projection of the 3{,}779 cosine-eligible responses
    (points, coloured by response language) with the eighteen
    Wikipedia anchors injected as stars (coloured by anchor
    language; labels merge anchors that land at the same
    position). Every English anchor sits in a separate
    English-response cluster, while the Russian and Ukrainian
    anchors for Bandera, Maidan, and MH17 land on top of each
    other inside mixed Russian--Ukrainian response regions.
    Visual adjacency in this projection and EVoC cluster
    membership are related but not identical: the MH17 anchor pair
    overlaps here yet sits in EVoC's noise partition, while the
    annexation pair --- visually separate --- shares a cluster by
    the co-clustering statistics in the text. The fused
    joint-cluster cases are Crimea, Maidan, and Bandera.}
  \label{fig:exp_016_anchor_map}
\end{figure}

\section{What cosine misses}\label{sec:exp_011}\label{sec:exp_003}\label{sec:mercury_failure}

The magnitude metric is silent on three categorically distinct
phenomena visible in the sweep.

YandexGPT returned the same Russian content-filter boilerplate on
every one of $270$ prompts across three languages, with no
successful response in any cell. The same model answers neutral
prompts in EN and RU without difficulty, so the refusal is
domain-specific. The cosine pipeline treats the response text as
null; the finding surfaces only because EVoC clustering produces a
tight pure-model cluster of identical refusal embeddings. This
recovers the embedding-space signature of structural censorship
inheritance that \citet{urman2025silence} predict but cannot verify
in their Western-only data.

Extending the Yandex sweep across the five conflict families of
\S\ref{sec:exp_006} --- a dedicated Yandex-only refusal sweep at
ten repeats per cell, giving $270$ calls per three-language family
and $180$ per two-language family --- shows the filter is graded
rather than uniform. Refusal varies steeply across families and is
highest where Russian state narratives are most directly
implicated: $100\%$ on the
imaginary Russo-Ukrainian event ($45/45$), $77\%$ on
Israel--Palestine ($208/270$), $52\%$ on India--Pakistan
($140/270$), $22\%$ on the Falklands ($40/180$), and $6\%$ on the
Taiwan strait ($10/180$). (The interest ordering is a descriptive
reading of the observed rates, not an ex-ante quantity; no
independent proxy for Russia-state interest is fitted.) Within events the filter is typed by
question form: the two one-sided-defence slots (q07, q08) are
filtered at ${\geq}50\%$ in every non-imaginary family, and at the
low-interest limit they are the \emph{only} filtered slots --- the
Falklands sweep filters both POV questions at $20/20$ and nothing
else, and Taiwan filters half of the pro-PRC POV slot and nothing
else. The filter is neither random nor domain-uniform; it is a
policy with a topology, and \S\ref{sec:yandex_limit} takes up what
that topology implies. Figure~\ref{fig:exp_022_cascade} shows both
views of the gradient.

\begin{figure}[!htbp]
  \centering
  \includegraphics[width=\linewidth]{exp_022_yandex_refusal_cascade.png}
  \caption[YandexGPT refusal cascade]{The YandexGPT content filter
    across five conflict families. Left: share of prompts refused
    per family, ordered by decreasing Russia-state interest, with
    refused/total counts inside the bars. Right: within-family
    refusal rate split into the two one-sided-defence slots
    (q07/q08, pooled) versus the remaining seven questions; at the
    low-interest limit (Falklands, Taiwan) only the POV slots are
    filtered. Bars use a refusal-intensity gradient rather than the
    family palette.}
  \label{fig:exp_022_cascade}
\end{figure}

Mercury-2, the US-trained diffusion LLM in the matrix, exhibits
language-conditional topic avoidance. Ten Mercury-2 responses per
language to the ``little green men'' question were coded against a
pre-specified (repository-frozen) four-level Crimea-coverage rubric (2~Primary /
1~Secondary / 0~Absent / R~Refusal) by \verb|claude-sonnet-4-6| at
temperature zero, with file order shuffled before coding; since the
coder reads the response text, the response language itself is
necessarily visible to it. The medians match the pre-specified
hypothesis exactly: EN $=1$ (Secondary), RU $=0$ (Absent), UK $=2$
(Primary). Ukrainian responses lead with the Crimea framing in
$9/10$ cases (the tenth is a refusal); Russian responses stay in
sci-fi or folklore framings in $7/10$ and reach Primary once, so
the avoidance is pronounced but not absolute. Same model, same
temperature, same question in three translations --- three
willingness levels to engage with the topic. The cosine number
averages this away.

The 7B-locals stylistic-fingerprint result from
\S\ref{sec:exp_016} is the third
category, reframed here as a scope condition: any
provider-agnostic claim in the chapter applies to frontier-scale
models only.

\section{Imaginary content separates direction from
  divergence}\label{sec:exp_004}\label{sec:exp_004_metric1}\label{sec:exp_004_metric2}\label{sec:two_shapes}

The cosine metric measures distance to a Wikipedia anchor the model
has plausibly read in training. The imaginary-conflict pilot
removes the anchor altogether by replacing the case study with a
fictional Russo-Ukrainian cross-border incident
(\emph{Tarasenko Bridge}, 14 March 2024, invented place names).
Five mirrored questions across the fourteen cosine-eligible
providers with three repeats per cell target $630$ completions;
$628$ return text. The two missing cells are Mercury-2 completions
that returned zero visible tokens with
\verb|finish_reason="length"| --- the diffusion architecture spent
its token budget on internal denoising --- and are retained as data
rather than re-rolled.

The pre-specified stance code (frozen in the repository before
scoring; judge
\verb|claude-sonnet-4-6|, disjoint from the swept-model set) scores
each response on a $-2/+2$ rubric whose sign convention is
\emph{positive = Russian-ward, negative = Ukrainian-ward}; the same
convention is used for every stance-judge number in this thesis
unless flagged otherwise. It codes $84\%$ of scored responses
neutral ($79\%$ of all completions; the difference is refusals,
which are excluded from scoring), with the residual tilting
\emph{Ukrainian-ward}: per-language means EN $-0.092$, RU $+0.010$,
UK $-0.170$. The originally-hypothesised ``models default to a
Russian framing on imaginary content'' is falsified. A second
frontier judge from a different provider reproduces the code's
directional signs on a stratified sample ($\kappa = 0.78$, no sign
inversions; \S\ref{sec:threat_single_judge}).

A within-model anchor-free divergence metric on the same response
set surfaces the sharper claim. Free-form prose has the
\emph{highest} cross-lingual content divergence in the bank
($2.86\times$ same-language baseline) but the \emph{lowest} stance
direction (mean $0.00$): models confabulate different things per
language without picking a side. Direct attribution has the
inverse: low content divergence ($1.73\times$) and the strongest
directional pull (UK mean $-0.56$, Ukrainian-ward). Bias on
imaginary content has
two distinct shapes --- \emph{content divergence without direction}
in free-form prose and \emph{directional commitment without content
divergence} in attribution. The literal national-stance hypothesis
collapses both into one and misses the structure. One provider
(\verb|command-r7b-arabic-02-2025|) shows the largest single-model
swing: a Russian-prompt mean of $+0.47$ (Russian-ward) against a
Ukrainian-prompt mean of $-0.60$ (Ukrainian-ward), a $1.07$-point
swing on the same fictional event.
Figure~\ref{fig:exp_004_two_shapes} plots the dissociation.

\begin{figure}[!htbp]
  \centering
  \includegraphics[width=0.85\linewidth]{exp_004_two_shapes.png}
  \caption[Two shapes of bias on imaginary content]{The five
    imaginary-conflict questions plotted by cross-lingual content
    divergence (as a multiple of the same-language baseline,
    Metric~2) against the largest per-language $|$stance
    mean$|$ (Metric~1). The free-form news-report question sits
    alone at maximal divergence with zero direction; the
    attribution questions (responsibility, belongs-to) sit at
    minimal divergence with the strongest direction. Content
    divergence and directional commitment dissociate by question
    form.}
  \label{fig:exp_004_two_shapes}
\end{figure}

\section{The pattern generalises across five
  conflicts}\label{sec:exp_006}\label{sec:exp_006_debiased}

The fourteen-provider sweep applies with no code changes to four
further conflict families: Israel--Palestine (EN/HE/AR),
India--Pakistan (EN/HI/UR), Taiwan strait (EN/ZH), Falklands
(EN/ES). Total spend $\$5{.}62$ across $9{,}023$ calls at $N{=}5$
samples per cell.

The English diagonal ranges from $+0.32$ on India--Pakistan to
$+0.05$ on the Falklands; in descending order, India--Pakistan
($+0.32$), Israel--Palestine ($+0.25$), Russo-Ukrainian ($+0.18$),
Taiwan ($+0.11$), Falklands ($+0.05$). These are raw values, and
the mechanical-ceiling correction of \S\ref{sec:exp_001} matters
across families even more than across embedders, because
anchor--anchor similarity varies sharply by family (English
ceilings: India--Pakistan $0.46$, Israel--Palestine $0.39$,
Russo-Ukrainian $0.24$, Taiwan $0.19$, Falklands $0.13$).
Ceiling-normalised, the ordering changes: the Russo-Ukrainian case
rises to the top ($73\%$ of ceiling attained), followed by
India--Pakistan ($70\%$), Israel--Palestine ($63\%$), Taiwan
($59\%$), and the Falklands ($37\%$), and the six-fold raw spread
compresses to under two-fold. Between-family magnitude comparisons
in this thesis are therefore reported raw but read through the
ceilings; the raw ranking partly reflects how distinct each
family's Wikipedia treatments are, not only how strongly models
pull toward them. Every English diagonal and
every contested-language diagonal is positive on raw cosine with
one exception: the Taiwan-strait Chinese cell sits marginally below
zero ($-0.008$, negative for seven of ten providers), the weakest
contested-language signal in the sweep. The Falklands near-null
functions as the falsification
anchor for the pipeline: the metric does not return a strong signal
wherever it is pointed. The natural explanation --- that the
Falklands War is the oldest conflict in this set, and
\citet{oeberst2020wikipedia}'s Study~2 documents a
recency--strength gradient --- is the one this chapter initially
adopted. \S\ref{sec:exp_007} withdraws it.

Cross-event debiasing produces an asymmetric pattern. The English
signal collapses by only $5\%$--$13\%$ in four of the five families
--- the central replication of the original Russo-Ukrainian
content-driven finding --- while the Falklands English signal, the
smallest in raw terms, collapses by $38\%$: where the raw pull is
weak, a larger share of it is language rather than content. The
contested-language signals collapse
asymmetrically: Hebrew by $96\%$, Ukrainian by $66\%$, Spanish by
$51\%$; Arabic ($20\%$), Hindi ($37\%$), and Urdu ($31\%$) collapse
only partially. (Collapse percentages are computed over the nine to
thirteen providers per family with both raw and debiased analyses
on disk, not the full fourteen.) The off-topic baseline of \S\ref{sec:topic_vs_lang}
converges on the same Hebrew-is-language-axis reading: Hebrew has
the lowest topic lift ($+0.18$) and the highest off-topic baseline
($0.36$ vs typical $0.10$--$0.16$).
Figure~\ref{fig:exp_006_dumbbell} shows the full raw-to-debiased
decomposition, including the Russo-Ukrainian case study of
\S\ref{sec:exp_001_debiased}, and
Figure~\ref{fig:exp_006_method_method} cross-plots the two
disentanglement methods directly: the cells where the language-axis
projection collapses the pull are the cells with the highest
same-language floor, with Hebrew the joint outlier on both
instruments. One uncontrolled alternative should be named: the
off-topic baseline is a mean cosine against randomly sampled
same-language control leads, and shorter, more generic leads have
mechanically higher mutual cosine, so a per-language baseline level
partly reflects how long that edition's control articles are; the
Hebrew reading is reported with that possibility open rather than
excluded.

\begin{figure}[!htbp]
  \centering
  \includegraphics[width=0.9\linewidth]{exp_006_method_vs_method.png}
  \caption[Two disentanglement methods cross-validated]{The two
    content-versus-language disentanglement methods cross-validated
    per (event, language) cell: language-axis debiasing collapse
    ($x$) against the off-topic same-language baseline ($y$), with
    dashed guides at the medians. Hebrew sits alone in the
    language-axis-dominated corner; the English cells cluster at
    low collapse and a low language floor. The Taiwan-strait
    Chinese cell is omitted because its raw pull is negative
    (\S\ref{sec:exp_006}).}
  \label{fig:exp_006_method_method}
\end{figure}

\begin{figure}[!htbp]
  \centering
  \includegraphics[width=0.92\linewidth]{exp_006_debias_dumbbell.png}
  \caption[Raw versus debiased pull across the five conflict
    families]{Raw (filled) versus language-axis-debiased (open)
    diagonal ingroup pull per (family, language), averaged over the
    providers with both analyses on disk; annotations give the
    collapse share. English diagonals (blue) survive debiasing
    with a $5\%$--$13\%$ collapse everywhere except the weak
    Falklands signal; contested-language diagonals (red) collapse
    asymmetrically, from Arabic's $20\%$ to Hebrew's near-total
    $96\%$. The Taiwan-strait Chinese cell is negative already in
    raw form.}
  \label{fig:exp_006_dumbbell}
\end{figure}

\section{A new metric demotes the flagship
  conflict}\label{sec:exp_021}

Cosine answers \emph{how topically similar is the response to its
own-language anchor?} It does not answer \emph{which side's framing
vocabulary does the response use?} The two have different answers
on the Second Intifada question for GPT-4o-mini: the Hebrew response
emphasises Hamas--PLO power struggles and reports Sharon's visit as
\emph{``perceived as a provocation''}; the Arabic response
emphasises Israeli settlement policy, ``Israeli forces' violent
response'', and the same visit \emph{``was a major provocation''}.
Row-centred cosine between the two against their anchors is
$+0.004$. The framing difference is invisible to the cosine number.

The paired-edit stance-axis projection of \S\ref{sec:stance_axis}
is built for exactly this question: a per-conflict axis fitted from
frozen minimal-edit sentence pairs, onto which each response
embedding is projected as a signed scalar.
Table~\ref{tab:exp_021_cross_conflict} reports the per-conflict
cross-language gap ranking;
Figure~\ref{fig:exp_021_cross_conflict} shows the per-provider
spreads. Under the cluster bootstrap of \S\ref{sec:stats}
(providers $\times$ responses) the mean gaps carry 95\% intervals
of $[0.107, 0.155]$ (India--Pakistan), $[0.103, 0.138]$
(Israel--Palestine), $[0.039, 0.078]$ (Taiwan), $[0.039, 0.065]$
(Falklands), and $[0.028, 0.043]$ (Russo-Ukrainian): the
Russo-Ukrainian floor is disjoint from the two leaders and touches
the Taiwan--Falklands pair, so the demotion of the flagship pair is
resolved against India--Pakistan and Israel--Palestine and
suggestive against the middle of the table.

\begin{table}[!htbp]
  \centering
  \begin{tabular}{lrrrr}
    \hline
    Conflict & Mean gap & Max single-provider gap
             & Axis seed sep. & $n$ providers \\
    \hline
    India--Pakistan        & $+0.131$ & $+0.173$ & $0.47$ & 13 \\
    Israel--Palestine      & $+0.117$ & $+0.222$ & $0.43$ & 13 \\
    Taiwan strait          & $+0.056$ & $+0.112$ & $0.52$ & 11 \\
    Falklands              & $+0.052$ & $+0.075$ & $0.62$ & 13 \\
    Russo-Ukrainian        & $+0.033$ & $+0.055$ & $0.51$ & 14 \\
    \hline
  \end{tabular}
  \caption[Stance-axis projection across five conflicts]{Per-conflict
    cross-language stance gap. The mean cross-language gap averages
    the per-provider max-minus-min stance score over the $n$
    providers with at least two language cells in the projection
    data (Yandex included where it answered); axis seed separation
    is the embedding distance between
    pole centroids before any response is projected. Bootstrap
    intervals for the mean gaps are given in the text.}
  \label{tab:exp_021_cross_conflict}
\end{table}

\begin{figure}[!htbp]
  \centering
  \includegraphics[width=\linewidth]{exp_021_cross_conflict_5panel.png}
  \caption[Stance-axis projection: five-conflict small
    multiples]{Paired-edit stance projection per provider across the
    five conflict families. The Russo-Ukrainian pair occupies the
    large left panel; the four further conflicts sit in the
    right-hand grid. Each row is one provider (sorted by its
    cross-language gap); dots are per-language cell means in the
    canonical language colours (per-panel legends), and the grey
    segment spans the provider's cross-language framing gap. Each
    panel's x-axis names its two framing poles; panel scales are
    independent, so spreads compare through the gap values of
    Table~\ref{tab:exp_021_cross_conflict}, not across panels by
    eye. Providers follow the same inclusion rule as the table.}
  \label{fig:exp_021_cross_conflict}
\end{figure}

The gap exists in every conflict. India--Pakistan tops the
ranking: Hindi and Urdu use different scripts and index two
distinct internal-narrative communities, and the gap is roughly
four times the size of the Russo-Ukrainian gap. The surprise is the
Russo-Ukrainian pair --- the conflict the project was built around,
and the one carrying the largest EN ingroup pull on cosine ---
showing the \emph{smallest} framing gap. The cosine and stance
rankings do not co-vary. Three non-exclusive readings are
consistent: providers may reproduce a single dominant
Russo-Ukrainian narrative across all three languages
(Western-framing dominance); Russian and Ukrainian may share enough
surface form that the embedder collapses the framing distinction
(Cyrillic-script convergence); or the intifada and Kashmir contests
may carry sharper community-internal narrative variation than
Crimea does. None is claimed on these data alone; each is a
testable extension. The between-family comparison also inherits the
two construction caveats of \S\ref{sec:stance_axis}: the axes are
fitted from English seed sentences, and the Taiwan and Falklands
gaps are ranges over two languages where the other families' are
ranges over three.

Freezing the lexicons before projection rules out tuning the axis
to the data, but not sensitivity to which of the twelve
hand-written pairs were chosen, so the
ranking is checked directly with a leave-one-pair-out ablation: for
each conflict every pair is dropped in turn and the axis re-fitted
from the remaining eleven, re-projecting the already-embedded
responses, for sixty re-fits in total and no new model calls. The
cross-language gap is stable under every deletion. The largest change
from removing any single pair is $0.007$, on Israel--Palestine, well
below the spacing that separates the clearly ranked conflicts, and
the five-conflict ordering holds across all sixty re-fits. The
margins are tight only between the adjacent Taiwan ($+0.056$) and
Falklands ($+0.052$) gaps, whose leave-one-out ranges nearly meet;
the India--Pakistan and Israel--Palestine lead and the
Russo-Ukrainian floor are unaffected. The ranking rests on the
lexicon as a whole, not on any individual pair. A separate
censoring check addresses Yandex, which enters the estimator only
where it answers and whose non-refused responses are therefore a
filtered sample: recomputing the gaps with Yandex excluded gives
$0.136$ / $0.108$ / $0.050$ / $0.055$ / $0.033$ --- the ranking is
unchanged except that the already-indistinguishable Taiwan and
Falklands pair swaps. Two further design-asymmetry checks close
off mechanical explanations for the Russo-Ukrainian floor: the
home sweep has $N{=}10$ samples per cell against $N{=}5$ elsewhere,
and a larger sample mechanically shrinks a max-minus-min range, so
the gap is recomputed with the home corpus subsampled to five per
cell (it stays at $+0.033$); and recomputing all five gaps on the
ten providers common to every family gives
$0.151$ / $0.107$ / $0.050$ / $0.057$ / $0.031$ --- floor and
leaders unchanged.

\section{The temporal gradient: eight conflict families across
  a millennium}\label{sec:exp_007}

The five-conflict scale-up of \S\ref{sec:exp_006} varies the
conflict but holds time roughly fixed: each family contributes a
single focal event. This leaves the second of
\citet{oeberst2020wikipedia}'s two structural moderators untested.
Their first moderator, editor homogeneity, has no direct analogue in
a language model. Their second is \emph{recency}: in the human
Wikipedia data, more recent conflicts carry stronger ingroup bias.
If a model reproduces the distribution of its training text on
contested events, as the source-side mechanism of
\S\ref{sec:mechanisms} supposes, it should inherit the recency
gradient along with the bias itself. Old events should read
neutrally; recent events should not.

That prediction is directly testable by holding the conflict family
fixed and varying the event's date. Each of eight families was given
an \emph{event ladder}: a set of focal events within the same
bilateral dispute, spanning as much of its documented history as
Wikipedia supports in both parties' languages. Six ladders cover
contested disputes --- Russia--Ukraine (nine events, 988--2022),
Israel--Palestine (1917--2023), India--Pakistan (1947--2019),
China--Japan (1894--2012), Greece--Turkey (1453--1974), and
Poland--Russia (1772--1944). Two further ladders are
single-event \emph{settled} baselines, chosen as bilateral
separations that concluded without lasting contestation: the Velvet
Divorce (1993) and the dissolution of the Norway--Sweden union
(1905). These supply a floor: whatever the metric reads on a
separation both parties regard as closed is the value that
bilingual difference alone produces, with no contestation left to
measure.

The ladder sweep covers 31 events in 16 languages. The cosine arm
aggregates the nine providers that completed the full cross-event
sweep --- Claude Haiku, DeepSeek, GPT-4o-mini, Gemini Flash, Jamba,
Mercury-2, Qwen-Plus, Grok-3-mini, and GLM-4.5; Baidu is excluded
by a dead API endpoint, the two Cohere models by trial-key quota,
the two locally hosted 7B models were not run on the ladder events,
and Yandex refuses categorically --- with nine-of-nine coverage on
every event except the 2022 invasion (eight of nine). The refusal
arm scans every ladder completion on disk, including providers
excluded from the cosine aggregation, for 31{,}457 scored
completions; the stance arm covers the post-1900 subset of the
ladder (23 of 31 events) by design. Like the embed-time detector
(\S\ref{sec:protocol}), the multilingual refusal classifier used
here is rule-based and not validated against human labels; a
stratified spot-check of flagged completions in the highest-rate
cells found genuine declines throughout, but on the
one-sided-defence slots many flagged completions decline the
\emph{requested framing task} while still engaging with the
facts --- ``refusal'' in this arm means refusal to comply with the
prompt, not silence.
One caveat governs the whole section and is repeated
where it bites: apart from English, and apart from the Russian and
Ukrainian prompts inherited from the case-study bank, every ladder
prompt is machine-translated and has not been reviewed by a native
speaker. The ladders are therefore exploratory in a way the
case-study sweep is not.

\subsection{Arm 1 --- cosine ingroup pull does not decay with
  time}\label{sec:exp_007_cosine}

Table~\ref{tab:exp_007_gradient} reports the English diagonal pull
across each ladder; Figure~\ref{fig:exp_007_overlay} overlays all
eight families on a common time axis.

\begin{table}[!htbp]
  \centering
  \begin{tabular}{lrrrr}
    \hline
    Family & Events & Span & EN pull range & Most recent \\
    \hline
    India--Pakistan   & 4 & 1947--2019 & $+0.251$--$0.290$ & $+0.251$ \\
    Israel--Palestine & 4 & 1917--2023 & $+0.202$--$0.266$ & $+0.225$ \\
    Greece--Turkey    & 4 & 1453--1974 & $+0.174$--$0.228$ & $+0.228$\rlap{$^\ast$} \\
    China--Japan      & 4 & 1894--2012 & $+0.091$--$0.209$ & $+0.177$ \\
    Russia--Ukraine   & 9 & \phantom{0}988--2022 & $+0.075$--$0.185$ & $+0.075$\rlap{$^\dagger$} \\
    Poland--Russia    & 4 & 1772--1944 & $+0.118$--$0.165$ & $+0.165$\rlap{$^\ast$} \\
    \hline
    \emph{Velvet Divorce (settled)}  & 1 & 1993 & $+0.070$ & --- \\
    \emph{Norway--Sweden (settled)}  & 1 & 1905 & $+0.092$ & --- \\
    \hline
  \end{tabular}
  \caption[Temporal gradient across eight conflict families]{English
    diagonal ingroup pull across each family's event ladder, ordered
    by ladder maximum. $^\ast$~marks families whose \emph{most
    recent} event is also their ladder maximum;
    $^\dagger$~marks the single collapse. The two settled baselines
    sit below the range of every contested family.}
  \label{tab:exp_007_gradient}
\end{table}

\begin{figure}[!htbp]
  \centering
  \includegraphics[width=\linewidth]{exp_007_conflict_overlay_enpull.png}
  \caption[Eight-family temporal overlay]{English diagonal ingroup
    pull for every ladder event, one slot per event in year order,
    in the canonical family colours. The grey band spans the two
    settled baselines (isolated $\times$ markers); the shaded right
    margin marks the live events (2022--2023). The spread of
    lines is vertical (between families) rather than horizontal
    (within families): a family's level is stable across its own
    history, and no family trends toward zero at its old end.}
  \label{fig:exp_007_overlay}
\end{figure}

Four findings follow.

\emph{Ingroup pull is positive on every one of the 31 events.} The
oldest event in the corpus, the 988 baptism of Kievan Rus', returns
$+0.126$ --- comparable to the 2014 case study that opened the
thesis ($+0.173$ on the nine-provider ladder subset; the
full-matrix figure of \S\ref{sec:exp_001} is $+0.175$) and well
above the settled floor. A thousand years
of elapsed time does not neutralise the effect.

\emph{The recency gradient does not transfer from human editors to
model responses.} Pooled across all 29 contested events, the rank
correlation between event date and English pull is
$\rho = {+}0.25$ --- weakly in the direction
\citet{oeberst2020wikipedia} predict. That pooled figure is
misleading. Computed \emph{within} families, where the comparison is
actually controlled, the correlation splits three against three:
Israel--Palestine $-0.80$, India--Pakistan $-0.40$, Russia--Ukraine
$-0.25$ against Poland--Russia $+0.40$, China--Japan $+0.20$,
Greece--Turkey $+0.20$. These are Spearman rank correlations over
small ladders --- four events per family, nine for Russia--Ukraine
--- and none is individually significant (the minimum attainable
two-sided $p$ at $n{=}4$ is $0.083$); what they jointly support is
the absence of a consistent gradient, not the presence of opposing
gradients within families. The pooled positive correlation is an
artefact of aggregation --- India--Pakistan contributes both the
most recent events and the highest baseline level, and pooling reads
that between-family level difference as a within-family time trend.
On the controlled comparison there is no consistent temporal
gradient in either direction. The model inherits the \emph{direction}
of Wikipedia's ingroup bias without inheriting its
\emph{recency structure}.

\emph{The settled baselines calibrate the metric --- with a
mechanical qualification the ceilings force.} In raw pull both
settled separations sit at $+0.070$ and $+0.092$, below the minimum
of every contested ladder except the single Russia--Ukraine 2022
cell (an exception that disappears under the length-matched anchor
control below). But the settled pairs are precisely the ones whose
languages are mutually intelligible, and mutual intelligibility is
what raises anchor--anchor cosine and lowers the mechanical
ceiling: the settled English ceilings are $0.17$ (Velvet Divorce)
and $0.16$ (Norway--Sweden) against $0.19$--$0.46$ for the
contested families. Ceiling-normalised, the Velvet Divorce attains
$40\%$ of its ceiling and Norway--Sweden $58\%$ --- the latter
above several contested ladder events and only marginally below
the Taiwan strait's $59\%$ from the five-conflict sweep. The
normalisation also has a consequence this paragraph should own
rather than leave for an examiner to find: the \emph{contested}
Falklands attains $37\%$, below both settled baselines. Under the
thesis's own correction, then, the metric cannot separate a
contested-but-weak dispute from no dispute at all --- which is
coherent with the corpus-pull reading of
\S\ref{sec:four_instruments} (the Falklands' English-language
treatment is close to the global English mainstream, so there is
little distinctive own-corpus signal to realise), but it forecloses
reading the settled floor as a clean contested/settled
discriminator. The raw settled-vs-contested separation is
therefore partly mechanical: the design choice that made the
baselines clean (closely related settled language pairs) also
compresses their attainable pull. What the baselines establish is
accordingly weaker than first appears: raw pull is low where the
two treatments are near-duplicates, whether the cause is
settlement or language proximity, and the settled floor cannot by
itself attribute the contested families' higher raw pull to
contestation. The two scales answer different questions and are
used accordingly: raw pull answers ``how much own-corpus
resemblance is there?'' (the settled-floor and matched-anchor
claims are raw-scale claims), while ceiling-normalised pull
answers ``how much of the attainable resemblance is realised?''
(the between-family comparison). The sign claim is unaffected
on either scale.

\emph{The live-event collapse is unique to Russia--Ukraine, which
falsifies a tempting general law.} The 2022 invasion returns
$+0.075$, less than half the family's ladder maximum and at the
settled floor. The obvious reading --- models hedge on live,
high-salience events --- does not survive contact with the other
ladders. Greece--Turkey peaks at Cyprus 1974 and Poland--Russia
peaks at the 1944 Warsaw Uprising, each family's \emph{most recent}
event being its \emph{strongest}, and Israel--Palestine's 2023 Gaza
cell sits mid-range rather than collapsed. Whatever suppresses the
2022 cell is specific to that event, not a property of recency. (The
Cyprus and Warsaw peaks do not survive the anchor-length control
below; the Gaza cell does, and India--Pakistan's 2019 cell becomes
its family's \emph{maximum} under matching --- so the
event-specificity argument stands on the controlled comparison too.)

This last finding forces a correction elsewhere in the thesis. The
Falklands near-null of \S\ref{sec:exp_006} was read as consistent
with the recency--strength gradient, the Falklands War being the
oldest of the five conflicts in that sweep. The ladder data do not
support that explanation: the 1453 fall of Constantinople returns
$+0.211$ and the 1772 partitions of Poland return $+0.143$, both far
above the Falklands' $+0.05$ despite being centuries older. Age
alone does not produce a weak signal. The Falklands result stands as
an observed near-null and continues to serve as a falsification
anchor for the pipeline, but the recency explanation for it is
withdrawn; what distinguishes the Falklands is more plausibly the
low contestation intensity of the dispute in present-day discourse,
which the ladder design does not isolate.

\paragraph{The anchor-size confound.} The most serious threat to the
Russia--Ukraine 2022 result is that its Wikipedia anchors are by a
wide margin the longest and most Russian-skewed in the corpus, so
event recency is confounded with anchor size. A longer anchor covers
more topical ground and can depress a cosine diagonal for reasons
that have nothing to do with framing.
Figure~\ref{fig:exp_007_anchor_control} plots pull against anchor
token length across the ladder. The relationship is weak enough that
anchor length does not account for the 2022 collapse on its own, but
it is not absent, and the 2022 cell is the highest-leverage point in
the plot. The matched-length re-run pre-specified as the remediation
was therefore carried out: every anchor lead was truncated to the
token length of the shortest anchor in its conflict family and
re-embedded (271 embedding calls, no new model calls), and the
row-centred diagonals were recomputed against the existing response
embeddings. Under matched anchors the 2022 cell rises from $+0.075$
to $+0.111$: taking the deficit as the gap below the family maximum
($0.185 - 0.075 = 0.110$), matching recovers $0.036$ of it ---
roughly a third of the deficit was anchor length. The
residual dip is real: the cell remains the Russia--Ukraine family
minimum (matched family range $+0.111$--$0.194$), though its margin
over the settled Velvet Divorce baseline is a point estimate of
only $+0.007$ --- and because matching truncates each anchor to its
\emph{family's} shortest lead, the matched cells being compared
across families are truncated to different lengths, so this
between-family margin inherits the anchor-geometry caveat of the
settled-baseline paragraph above. The 2022 cell is also the only ladder cell with
eight- rather than nine-provider coverage; recomputing the family
on the common eight-provider subset moves every other cell by at
most $0.003$ and leaves 2022 the family minimum, so provider
composition does not explain the dip. The
within-family recency correlations are also unstable under matching
(three of six flip sign), which reinforces the conclusion that no
consistent temporal gradient exists in either anchor regime.
Residual caveats are noted in \S\ref{sec:threat_anchor_length}.

\begin{figure}[!htbp]
  \centering
  \includegraphics[width=0.92\linewidth]{exp_007_fig6_anchor_length_control.png}
  \caption[Anchor-length control]{English diagonal pull against
    the mean anchor lead length (in tokens, across the event's
    languages) for every ladder event, in family colours; settled
    baselines as $\times$ markers. No strong relationship is
    visible, but the Russia--Ukraine 2022 cell combines the longest
    contested anchors with the lowest pull, making it the
    highest-leverage point for the confound between event recency
    and anchor size.}
  \label{fig:exp_007_anchor_control}
\end{figure}

\subsection{Arm 2 --- stance projection reproduces the
  Russo-Ukrainian floor}\label{sec:exp_007_stance}

The paired-edit stance axis (\S\ref{sec:stance_axis}) was extended
to the ladder corpus, with a lexicon per family, over the
post-1900 event subset. Because each lexicon
fixes its own pole assignment, the absolute stance level is not
comparable across families; only the cross-language \emph{gap} is.
Table~\ref{tab:exp_007_stance} reports the gaps; they are
descriptive point estimates (\S\ref{sec:stats}), so adjacent
entries separated by a few thousandths are ties at this
precision.

\begin{table}[!htbp]
  \centering
  \begin{tabular}{lr}
    \hline
    Family & Cross-language stance gap \\
    \hline
    India--Pakistan   & $0.142$ \\
    Greece--Turkey    & $0.123$ \\
    Israel--Palestine & $0.097$ \\
    Poland--Russia    & $0.086$ \\
    \emph{Norway--Sweden (settled)} & \emph{$0.053$} \\
    China--Japan      & $0.036$ \\
    Russia--Ukraine   & $0.033$ \\
    \emph{Velvet Divorce (settled)} & \emph{$0.028$} \\
    \hline
  \end{tabular}
  \caption[Stance gap across the eight ladders]{Cross-language stance
    gap per family, computed as the spread between the highest and
    lowest per-language mean projection. Settled baselines are
    italicised; note that one of them does \emph{not} sit at the
    floor.}
  \label{tab:exp_007_stance}
\end{table}

Two results and one caveat.

The Russia--Ukraine gap is $0.033$, close to the five-conflict
stance sweep's value (\S\ref{sec:exp_021}). This agreement is a
\emph{stability} check, not convergent validity from an independent
method: the ladder arm imports the same frozen lexicon, the same
axis construction, the same projection code, and the same embedder,
and one of its five Russia--Ukraine events \emph{is} the
\texttt{ru\_uk\_core} sweep corpus itself. Excluding that shared
event, the remaining four historical events give a gap near
$0.04$ ($0.039$ pooled across events; $0.043$ as a mean of
per-event gaps) --- the home-case floor is stable when four further
events are added under the same axis, which is the claim the
agreement supports.
Figure~\ref{fig:exp_007_convergent} shows the two arms against each
other. India--Pakistan moves from $+0.131$ to $0.142$ and
Israel--Palestine from $+0.117$ to $0.097$ between the two sweeps,
so cross-arm agreement is loosest exactly where the event sets
differ most.

The Russo-Ukrainian floor from \S\ref{sec:exp_021} therefore
survives a second, larger corpus. The conflict the project was built
around continues to show the smallest cross-language framing gap of
any contested family measured, now against eight families rather
than five.

\emph{The caveat is that the settled/contested separation is far
less clean here than on cosine --- and the family-level numbers
carry real event-to-event spread.} The per-family figures are means
over unequal event counts (five for Russia--Ukraine, four or three
for most contested families, and a \emph{single} event for each
settled baseline), with no interval attached; the per-event gaps
show the spread. Russia--Ukraine's five events range from $0.014$
to $0.086$ and Greece--Turkey's two from $0.081$ to $0.166$, so
Norway--Sweden's single-event $0.053$ --- above China--Japan
($0.036$) and Russia--Ukraine ($0.033$) on the family means ---
falls \emph{inside} the Russia--Ukraine event-to-event range. On
the cosine arm every settled baseline sat below every contested
family in raw pull; on the stance arm one settled point estimate
does not, but with one event and zero degrees of freedom it cannot
be distinguished from the contested families' internal spread
either. The honest statement is that the stance axis is a weaker
discriminator between contested and settled disputes than anchor
cosine, and the Russo-Ukrainian floor is better described as ``at
the level of a settled dispute'' than as ``distinguishable from
one''.

\begin{figure}[!htbp]
  \centering
  \includegraphics[width=0.92\linewidth]{exp_007_convergent_validity.png}
  \caption[Convergent validity across stance sweeps]{Cross-language
    stance gap measured on the five-conflict sweep against the same
    quantity measured on the eight-family ladder corpus, for the
    families common to both. Russia--Ukraine lands on the identity
    line; the other shared families do not.}
  \label{fig:exp_007_convergent}
\end{figure}

\subsection{Arm 3 --- refusal is near-null, and that is the
  point}\label{sec:exp_007_refusal}

Every ladder completion was scored for refusal against the
classifier used in the Yandex sweep of \S\ref{sec:exp_011}. If the
2022 cosine collapse were produced by models declining to answer,
the refusal rate on that cell would be conspicuous.

It is not. Per-family mean refusal spans $0.58\%$ (Greece--Turkey)
to $2.33\%$ (China--Japan), with Russia--Ukraine at $1.41\%$ and
the two settled baselines lowest of all ($0.27\%$ Norway--Sweden,
$0.00\%$ Velvet Divorce). The family mean is the wrong denominator
for the one cell the question targets, so the check is also run at
cell level: within the 2022 cell the classifier flags $18.5\%$ of
DeepSeek's completions and $3.7\%$ of GPT-4o-mini's (soft refusals
that remain \emph{in} the cosine pool, so the dip cannot come from
dropping them), while English-language refusal in that cell is
$0.2\%$ --- the language whose diagonal defines the collapse is
essentially fully answered. What the classifier cannot see is
answered-but-deflecting prose, which would depress a diagonal
while counting as engagement; the arm closes off silence, not soft
deflection, and the anchor-length control carries the rest of the
explanation. The mechanism behind the collapse is what
the models say, not whether they say it ---
Figure~\ref{fig:exp_007_refusal} shows the full distribution.
This arm is a near-null on hard refusal, and its value is
precisely that: it closes
off the silence explanation for the headline result.

\begin{figure}[!htbp]
  \centering
  \includegraphics[width=0.92\linewidth]{exp_007_refusal_by_conflict.png}
  \caption[Refusal rate by conflict family]{Refusal rate per
    (event, language) cell, grouped by conflict family (canonical
    family colours; black dashes mark family means; cells above
    five per cent are labelled). Rates sit
    between one and two per cent almost everywhere; the settled
    baselines are lowest. The 2014 case-study cells clear five per
    cent in all three languages only because that corpus includes
    YandexGPT's categorical refusals; the remaining outliers are
    the genuinely language-conditional cells named in the text.}
  \label{fig:exp_007_refusal}
\end{figure}

The interesting structure is in the tail, though part of the tail
is an artefact. The 2014 case-study cells read high in all three
languages ($5.8\%$ EN, $6.9\%$ RU, $6.5\%$ UK) only because that
corpus contains YandexGPT's 270 categorical refusals; excluding
Yandex they fall to $0.2\%$, $1.5\%$, and $1.2\%$ respectively. The
genuinely language-conditional cells are the 1937 Nanjing Massacre
asked in Japanese ($7.9\%$) and Chinese ($5.8\%$), the 1940 Katyn
massacre asked in Russian ($6.2\%$), and the 2022 invasion asked in
Russian ($5.9\%$) --- each a case where the question is
domestically sensitive in the language it is being asked in. At
provider level the tail is dominated by DeepSeek, which refuses
$39.5\%$ of Nanjing-1937 completions, $23.5\%$ of Katyn, and
$18.5\%$ of the 2022 invasion. The
Nanjing cell is the quantified form of a qualitative pattern
visible in the raw responses and shown in
Figure~\ref{fig:exp_007_nanjing}: asked about the massacre in
Japanese, models are most likely to decline; asked in Chinese, they
answer with the framing standard in Chinese-language sources; asked
in English, they answer neutrally. Refusal is rare overall, but
where it occurs it is language-conditional in the same direction as
the framing effect this thesis measures --- which makes it a small
confirming instance of the main finding rather than a competing
explanation for it.

\begin{figure}[!htbp]
  \centering
  \includegraphics[width=\linewidth]{exp_007_qualitative_nanjing.png}
  \caption[Language-conditional response to the Nanjing
    Massacre]{Responses to the same question about the 1937 Nanjing
    Massacre in Japanese, Chinese, and English. The Japanese cell
    carries the highest refusal rate of any cell in the ladder
    corpus ($7.9\%$); the Chinese and English cells answer, with
    visibly different framing. Card accent colours are decorative
    and not part of the language palette.}
  \label{fig:exp_007_nanjing}
\end{figure}

\subsection{Arm 4 --- an LLM judge, read as
  corroboration only}\label{sec:exp_007_judge}

The fourth arm applies the LLM-as-judge stance protocol of
\S\ref{sec:anchor_free} to the most recent event in each family.
The judge is again \verb|claude-sonnet-4-6|; the judged responses
come from two providers (DeepSeek and GPT-4o-mini) over three
questions, for 144 judgments in total. The rubric is a per-family
$-1/{+}1$ axis whose positive pole is the second-named party of the
family --- for Russia--Ukraine, positive is \emph{Ukrainian}-ward,
the reverse of the $-2/{+}2$ convention of \S\ref{sec:exp_004}; all
signs in this subsection follow the per-family convention. At six
judgments per (family, language) cell this arm is far smaller than
the other three and is reported as corroboration, never as a
headline.

Russia--Ukraine is the only contested family in which every
language mean is non-negative: English $+0.225$, Russian $+0.025$,
Ukrainian $+0.408$ (Ukrainian-ward positive). Every other
contested family returns at least one negative language mean, and
India--Pakistan ($-0.49$ to $-0.63$) and China--Japan ($-0.25$ to
$-0.57$) are negative throughout
(Figure~\ref{fig:exp_007_judge}). The Ukrainian-ward tilt on
Russo-Ukrainian content is thus reproduced on a second corpus and a
second pair of judged providers --- but by the \emph{same} judge
model that produced the \S\ref{sec:exp_004} stance code, so the
agreement is corroboration across corpora, not across judges; a
latent judge-side preference would survive it intact
(\S\ref{sec:threat_single_judge}).

\begin{figure}[!htbp]
  \centering
  \includegraphics[width=0.92\linewidth]{exp_007_judge_stance_by_conflict.png}
  \caption[Judge stance by conflict family]{Mean LLM-judge stance
    score per (family, language) cell. Russia--Ukraine is the only
    contested family with no negative language mean. The settled
    Norway--Sweden baseline scores strongly negative, which a
    settled dispute should not, and is the basis for treating this
    arm as exploratory.}
  \label{fig:exp_007_judge}
\end{figure}

The reason for the caution is visible in
Figure~\ref{fig:exp_007_judge}. The
settled Norway--Sweden baseline returns means between $-0.25$ and
$-0.43$, a directional commitment on a dispute that has none to
make. Either the judges are reading ordinary historical narration as
partisan, or the rubric does not transfer to pairs without a live
contest. The other settled baseline, the Velvet Divorce, behaves as
expected ($-0.05$ to $+0.07$), so the failure is not uniform. Until
a judge protocol is validated on settled pairs, this arm can support
a direction that other instruments have already found; it cannot
establish one alone.

\subsection{What the four arms establish
  together}\label{sec:exp_007_synthesis}

The ladder corpus makes the headline finding harder to dismiss and
narrower than it first appeared. Ingroup pull is present on every
event tested across a millennium and eight bilateral disputes, it
separates from two settled-dispute baselines (narrowly at the 2022
cell, where the matched-anchor margin is $+0.007$), and it is not
produced by refusal. Against that, the recency gradient that the
Wikipedia literature would predict does not appear under controlled
within-family comparison, one prominent explanation offered earlier
in this chapter for the Falklands near-null is withdrawn as a
result, the Russo-Ukrainian 2022 collapse recovers about a third of
its deficit under
length-matched anchors without vanishing, having been confounded with
anchor size, and the stance and judge arms each fail a
settled-baseline sanity check that the cosine arm passes. The
strongest statement the data support is that the effect is real,
broad, and era-insensitive --- not that it is uniform, and not that
its magnitude is well understood.

The reconciliation panel that closes this chapter is a different
quartet from the four ladder \emph{arms} above: it keeps cosine,
stance, and judge, drops the refusal arm (a near-null here), and
adds the published forced-choice result as an external anchor.
These four measurement instruments now disagree on the same
Russo-Ukrainian conflict. Anchor cosine pulls English-ward on the
real-event corpus (\S\ref{sec:exp_001}, $+0.175$); stance
projection compresses toward zero on that same corpus ($+0.033$);
the Sonnet stance code on the \emph{imaginary} corpus
tilts Ukrainian-ward (\S\ref{sec:exp_004_metric1}, mean $-0.084$);
and \citet{li2024borderlines}'s discrete MCQ measurement on the
same Crimea item --- an external published result on vanilla GPT-4,
a model outside this matrix, not a measurement of this corpus ---
flips Russian-ward. The reconciliation
is the spine of Chapter~\ref{ch:discussion}.
